\section{Related work}
\label{sec:related}

The v0.4 study surveyed inference, search, team equilibria and signalling, and
that survey stands. This section covers what the present report additionally
depends on, and it is organised around a gap: a coverage sweep of the eight
reviews assembled for v0.4 finds no occurrence of \emph{war of attrition},
\emph{minimax regret}, \emph{distributionally robust}, \emph{Dynkin},
\emph{restricted Nash}, \emph{data-biased} or \emph{safe opponent exploitation},
and one occurrence of \emph{deadlock}. Those are exactly the two bodies of work
this report needs: the theory of voluntary stopping when both sides may wait, and
the theory of performing well against the worst opponent style rather than the
average one.

\subsection{Fish and Literature}

The published material on this game remains human and informal. McLeod's Pagat
entry \cite{pagat} is the canonical rules source and records the points on which
dialects differ; Develin's chapter \cite{develin} is the most detailed strategy
writing located, and supplies two observations this design uses --- that an ask
teaches the opponents as much as it teaches a teammate, and that pre-agreed
conventions are outside the game. A repeat search on 22 August 2026 --- the arXiv
API for \texttt{"Canadian Fish"}, \texttt{"Literature card game"} and
\texttt{"half-suit"}, returning zero results, plus a web sweep that surfaced only
adjacent games --- reproduced the v0.4 finding. Every novelty statement in this
report therefore takes the form ``we found no prior published work that \dots''
and should be read against that scope. The one agent located, Somani's
\texttt{literature} \cite{somani}, plays the four-player 48-card variant and
publishes no strength numbers; a closed-source mobile application advertises
six-player bots with no stated methodology \cite{playstore}; Literature is not
among the environments packaged with RLCard \cite{rlcard}.

\subsection{Voluntary stopping when waiting is free}
\label{sec:related-stopping}

The declaration decision in Fish is a multi-player, nonzero-sum, undiscounted
optimal-stopping game: each side chooses when to claim, the payoff depends on who
claims first and whether the claim was right, and waiting costs nothing in the
rules. Games of this shape are known to be badly behaved. Christensen,
Lindensj\"o and Neumann \cite{christensen-dynkin} record that for undiscounted
Dynkin games only $\epsilon$-equilibria are established in general, with
counterexamples showing that Nash equilibria need not exist even for
deterministic stationary rewards, and that this survives enlargement to
randomised stopping times. Hamad\`ene, Hassani and Morlais
\cite{hamadene-dynkin} prove existence of an $\epsilon$-Nash equilibrium for the
$N$-player discrete-time case by a constructive algorithm, under a payoff
structure that is the Fish structure: payoffs depend on the set of players who
stop at the first stage at which anyone stops. These are results for general
Markovian reward processes rather than theorems about Fish, so the correspondence
is one of structure and is used here to name the missing ingredient, not to
assert a property of the game.

The classical model of mutual refusal to move is the war of attrition. Hendricks,
Weiss and Wilson \cite{hendricks-attrition} characterise the complete-information
continuous-time case and find a continuum of equilibria including arbitrarily
long delay: delay is an equilibrium phenomenon that a symmetric deterministic
pair will select, and \vfour{}'s mirror match is exactly symmetric and exactly
deterministic. The constructive converse is Mguni \cite{mguni-actioncost}: impose
a strictly positive cost on each action and the stochastic game acquires a unique
value and a Markovian saddle point at strategically chosen stopping times. That
is the cleanest published statement of the natural fix, and \S\ref{sec:diagnosis}
reports that in this game it does not work --- the cost is a mean-shifter that
leaves the tail of the distribution untouched.

There is a reason to be careful with holding costs beyond their measured
ineffectiveness. Admati and Perry \cite{admati-perry} show that delay is used as
a signalling device precisely because it is a more efficient signal of strength
than the alternatives, and that the delay does not vanish as the interval between
moves goes to zero. Applied to Fish this is Farrell--Gibbons two-audience cheap
talk \cite{farrell-gibbons} on the \emph{timing} channel: if how long a seat
waits before declaring depends on its private confidence, three opponents can
read that confidence. Any holding cost must therefore be a deterministic function
of public state, or be deliberately randomised.

We found no substantive literature on self-inflicted non-termination in
game-playing agents. Three distinct query formulations returned only anecdote,
resolved in every case by an arbitrary move cap --- which is what \vfour{} does,
and what \S\ref{sec:diagnosis} measures the cost of. That absence is itself a
result: voluntary-claiming games are an under-studied class and Fish is a clean
instance of one.

\subsection{Information-seeking actions, and what an ask leaks}
\label{sec:related-info}

An ask in Fish is simultaneously a request for material, a certificate to one's
teammates and a disclosure to one's opponents. The v0.4 review formalised the
trade-off as a per-decision secrecy rate in the Wyner \cite{wyner} sense --- the
mutual information the action carries to teammates minus the mutual information
it carries to opponents --- and built none of the four instruments it listed.
\vfour{} contains three information terms in its ask score and all three are
negative; \S\ref{sec:diag-features} gives the arithmetic.

Two results bear directly on the repair. Alvim et al.\ \cite{alvim-leakage} model
information flow as a game in which the defender's utility is a convex function
of the \emph{distribution} over its actions rather than the expected utility of
the action played, and show that a Nash equilibrium nevertheless exists. The
consequence for Fish is sharp: a deterministic argmax ask policy has maximal
leakage by construction, so no reweighting of leak-penalty features can reduce
leakage below the deterministic bound. Reducing it requires randomising over
near-tied asks. The second consequence is a caveat --- because utility is convex
in the mixed strategy, the leakage game is not bilinear and standard regret
matching does not apply unmodified.

Lervold, Peterson and King \cite{lervold-infogain} appear to be the only
published work that places an explicit information-gain incentive inside an
information-set tree search, reporting improvement over baseline and over a
cheating search in a hidden-information multi-action wargame. We were unable to
obtain the full text, so the finding is used here and the formula is not.
Serrino et al.\ \cite{serrino-deeprole} integrate deductive reasoning directly
into vector-form counterfactual regret minimisation for five-player Avalon; it is
the closest published system to Fish's structure, a hidden-role adversarial team
game in which every action is public and its informational content is the whole
game, and it is the existence proof that hard deduction can live inside a solver
rather than beside it.

\subsection{Online opponent modelling, and robustness to being modelled}
\label{sec:related-opponent}

\vfour{} models the other five seats with two global scalars, identical for every
opponent and never updated during a game (\S\ref{sec:diag-prior}). The literature
on doing better within a single short episode is mature and was absent from the
v0.4 corpus. Ganzfried, Wang and Chiswick \cite{ganzfried-multiplayer} give the
only opponent-modelling result located that is explicitly in the more-than-two
player regime. Rosman, Hawasly and Ramamoorthy \cite{rosman-bpr} pose exactly
this problem --- select online, from a library of policies, a response to a novel
task instance, under a short episode --- and a Fish game is
\numMirrorEvents{} public events long. Smith et al.\ \cite{smith-qmixing} learn a
value function against each pure opponent and combine them for any mixture
without retraining, with a classifier that refines the mixture weights during the
episode.

The robustness half matters more here, because the behaviour that prompted this
work is an opponent deliberately feeding a model. Johanson, Zinkevich and Bowling
\cite{johanson-rnash} give the shape of the dial: solve a game in which the
opponent plays a fixed model with probability $p$ and is free otherwise, so that
$p=1$ is a best response and $p=0$ an equilibrium, and intermediate $p$ traces an
exploitation/robustness frontier. Johanson and Bowling \cite{johanson-dbr} set
$p$ per information set as a function of how much data is actually available
there --- the right refinement for Fish, where the per-half-suit evidence about a
given opponent is a handful of observations, and the one that makes a
deliberately silent opponent fall back to the population prior instead of being
misread. Ganzfried and Sandholm \cite{ganzfried-safe} guarantee at least the game
value per period while still exploiting mistakes; M\"uller et al.\
\cite{muller-bestofboth} give the modern bandit-feedback formalisation of the
same idea. Yang et al.\ \cite{yang-bayestomop} handle opponents that are
themselves running Bayesian opponent models, and Milec, Kova\v{r}\'ik and Lis\'y
\cite{milec-abd} adapt to sub-rational opponents beyond a depth limit while
remaining robust against rational play.

\subsection{Worst case rather than mean case}
\label{sec:related-worstcase}

The standing evaluation preference of this project --- report the per-opponent
breakdown and an explicit worst case, never an aggregate alone --- has a matching
equilibrium concept. Aggarwal and How \cite{aggarwal-prmre} define a
probabilistically robust minimax-regret equilibrium for adversarial team games
under asymmetric information, minimising worst-case regret over a
high-confidence subset of the type space rather than over the whole of it, and
solve it inside a robust double-oracle loop. Translated to Fish, the type is the
opponent style, the high-confidence subset is the style set actually built, and
the prescription is to optimise worst-case regret rather than either the mean or
the absolute worst case --- promoting a reporting rule to a training objective.
Two practical notes come with it. Xu et al.\ \cite{xu-mirror} observe that
computing max-regret requires the per-style optimum first, which in this setting
is one arena run per style, so minimax regret is available as a post-hoc
selection criterion over checkpoints and not only as a training objective.
Beukman et al.\ \cite{beukman-remidi} document the failure mode of a naive
worst-case objective: once the bound is attained everywhere, learning stagnates
because the objective stops distinguishing policies --- which is a direct warning
against a raw minimum over a small style set. Vinitsky et al.\
\cite{vinitsky-rap} report that a single adversary produces exploitable policies
and replace it with a sampled population.

\subsection{Team games, and the concepts this report does not compute}

Three teammates share a payoff, hold different information and cannot
communicate, so the target concept is a team-maxmin equilibrium with correlation
\cite{celli-gatti}. Nothing here computes one. Two recent results scope the
ambition. Anagnostides et al.\ \cite{anagnostides-teamzero} show that computing
Nash equilibria in team zero-sum games is PPAD-complete, and remains so at
inverse-polynomial precision and with two-player teams; that concerns the
uncorrelated equilibrium, which is what three independently acting seats without
a shared per-episode seed actually implement. Liu et al.\ \cite{liu-hpsro}
identify a concrete defect in Team-PSRO \cite{mcaleer-teampsro}: its policy
sharing cannot cover the joint policy space in \emph{heterogeneous} team games
where teammates play distinct roles, yielding higher exploitability. Fish seats
are heterogeneous in the relevant sense, since turn flow and distance to the next
opponent differ by seat, so any future population loop for this game should be
the heterogeneous variant.

Two design questions in \S\ref{sec:mechanisms} are the Hanabi-style zero-shot
coordination questions \cite{bard-hanabi} in another game: how much a policy may
rely on its partner running the same blueprint, and how to keep a convention from
degenerating into a private code. Off-belief learning \cite{hu-obl} and
other-play \cite{hu-otherplay} are the reference treatments, and the level of
grounding is the natural explicit knob for the partner-regime distinction of
\S\ref{sec:mech-partner}.

\subsection{Determinization, and why it is not the foundation used here}

Perfect Information Monte Carlo --- sample hidden cards consistently with the
public history, solve the resulting perfect-information game, average
\cite{ginsberg-gib} --- is unsound however many worlds are drawn
\cite{frank-basin-ai}, and is degenerate in Fish for a structural reason as well:
a clairvoyant player never fails an ask, so almost every reachable position
carries the same perfect-information value and the average collapses to choosing
the ask most likely to hit. The v0.4 study gives the collapse in full. It is
worth restating here because the collapse target --- ``choose the ask most likely
to hit'' --- is a policy that would not have deadlocked. \vfour{}'s pathology is
not that it is a hit-maximiser; it is that its ownership features outbid the
hit-probability term at $p=0$ (\S\ref{sec:diag-features}). The policy-aware
reweighting of a rules posterior that the Skat literature uses
\cite{solinas-supervised,rebstock-inference,buro-kermit} is the correct frame for
the prior that \S\ref{sec:diag-prior} examines, and the counting machinery
underneath \cite{lsw,cryan-dyer} is unchanged from v0.4.
